\subsection{内部坐标束屏幕的原子可积性、Poisson--Gumbel 临界窗与电阻--Shapley 同一}\label{subsec:conclusion-coordinatebundle-integrable-poisson-gumbel-resistance-shapley}

沿用小节 \ref{subsec:conclusion-screen-hypercube-boundary-torsor-shapley}、
小节 \ref{subsec:conclusion-coordinatebundle-kirchhoff-poissonbinomial-logscale} 与
小节 \ref{subsec:conclusion-coordinatebundle-codimension-tree-star-kirchhoff-rigidity}
的记号。固定非空
\[
J\subseteq [n],
\qquad
s:=|J|,
\qquad
L_J:=2^{m(n-s)},
\qquad
S_J^{\mathrm{int}}\subseteq E(\Gamma),
\]
并记
\[
\mathcal D_{m,n}:=\{2^{m\ell}:0\le \ell\le n\}.
\]
对每个
\[
a_{J^c}\in\{0,\dots,2^m-1\}^{J^c},
\]
记对应板层为 \(\Sigma_{a_{J^c}}\)，其外边界面集为
\[
E_{a_{J^c}}:=\partial_{\mathrm{out}}\Sigma_{a_{J^c}},
\qquad
d_{m,s}:=|E_{a_{J^c}}|=2s\,2^{m(s-1)}.
\]
前述三条链已经分别固定了：后验纤维的仿射超立方体结构、外边界精确化的逐板层完全因子化、随机外边界审计的 Poisson--binomial 残余律、面级 Shapley--Banzhaf 原子律，以及严格边界价差下的纯外边界星树最优性。把这些接口并置之后，内部坐标束屏幕还可进一步压缩成一条更强的可积主线：原始隐藏复杂度在唯一无交互的对数标度下退化为纯 \(m\)-比特坐标原子；均匀外边界审计具有精确的 Poisson--Gumbel 临界窗；同层 Gibbs 边占用、有效电阻与 Shapley 价值在无权情形下严格重合；而严格边界价差又把零温最优补全压缩为纯星子族，并给出每个隐藏原子的显式熵密度。

\begin{theorem}[隐藏复杂度的原子可积化与对数标度唯一性]\label{thm:conclusion-coordinatebundle-atomic-integrability-logscale}
定义隐藏体积
\[
\Delta_m(J):=L_J=2^{m(n-s)}.
\]
则下列四类复杂度在内部坐标束屏幕上完全锁定于同一原始尺度，其中 Kolmogorov 项只差 \(O(1)\)：
\[
H(X\mid Y_J),
\qquad
\min\{\text{exact audit bits}\},
\qquad
\min\{\text{linear completion dimension}\},
\qquad
\max_{y}\max_{x:T_Jx=y}K(x\mid y,m,n,J).
\]
更精确地，前三者都恰等于
\[
L_J,
\]
而第四者满足
\[
\max_{y}\max_{x:T_Jx=y}K(x\mid y,m,n,J)=L_J+O(1).
\]

进一步，对任意函数 \(\phi:\mathcal D_{m,n}\to\RR\)，定义
\[
h_\phi(J):=\phi(\Delta_m(J)).
\]
则“所有二阶及以上 Möbius 原子同时消失”当且仅当存在常数 \(A_0,C\in\RR\) 使
\[
\phi(t)=A_0+C\log_2 t
\qquad
(t\in \mathcal D_{m,n}).
\]
因此一切无交互重标都只能写成
\[
h_\phi(J)=A_0+Cm(n-s),
\]
并对任意 \(j\notin J\) 满足
\[
h_\phi(J)-h_\phi(J\cup\{j\})=Cm.
\]
特别地，在自然规一化 \(h_{\log}(J):=\log_2\Delta_m(J)\) 下，隐藏复杂度严格分裂为 \(n-s\) 个彼此独立的 \(m\)-比特坐标原子。
\end{theorem}
\begin{proof}
由定理 \ref{thm:conclusion-coordinate-bundle-fourfold-complexity-collapse}，
\[
H(X\mid Y_J)=L_J,
\qquad
\min\{\text{exact audit bits}\}=L_J,
\qquad
\min\{\text{linear completion dimension}\}=L_J,
\]
并且
\[
\max_{y}\max_{x:T_Jx=y}K(x\mid y,m,n,J)=L_J+O(1).
\]
这就给出原始尺度上的四重一致坍缩。

另一方面，定理 \ref{thm:conclusion-coordinate-bundle-logscale-unique-nointeraction} 已证明：对
\[
h_\phi(J)=\phi(\Delta_m(J)),
\]
其全部二阶及以上 Möbius 系数同时为零，当且仅当 \(\phi\) 在 \(\mathcal D_{m,n}\) 上是 \(\log_2 t\) 的仿射函数。于是
\[
h_\phi(J)=A_0+C\log_2\Delta_m(J)=A_0+Cm(n-s),
\]
故对每个 \(j\notin J\) 都有
\[
h_\phi(J)-h_\phi(J\cup\{j\})=Cm.
\]
取 \(C=1\) 即得自然的 \(\log_2\) 标度下的 \(m\)-比特原子分裂。证毕。
\end{proof}

\begin{theorem}[均匀外边界审计的 Poisson--Gumbel 临界窗]\label{thm:conclusion-coordinatebundle-poisson-gumbel-critical-window}
设每张外边界面独立地以相同概率 \(p\in[0,1]\) 被纳入审计集 \(B\)。记随机残余歧义秩为
\[
\Lambda_J(B):=\rank_{\ZZ}\ker f_{S_J^{\mathrm{int}}\cup B}.
\]
则
\[
\Lambda_J(B)\sim \mathrm{Bin}\!\bigl(L_J,q_{m,s}(p)\bigr),
\qquad
q_{m,s}(p):=(1-p)^{d_{m,s}}.
\]
从而 exactification 概率具有闭式
\[
\mathbb P\bigl(\Lambda_J(B)=0\bigr)=\bigl(1-q_{m,s}(p)\bigr)^{L_J}.
\]

对任意 \(c\in\RR\)，定义临界窗
\[
p_{m,s}(c):=1-\exp\!\left(-\frac{\log L_J+c}{d_{m,s}}\right).
\]
则有精确恒等式
\[
q_{m,s}\bigl(p_{m,s}(c)\bigr)=\frac{e^{-c}}{L_J},
\]
并且
\[
\Lambda_J(B)\ \Longrightarrow\ \mathrm{Poisson}(e^{-c}),
\qquad
\mathbb P\bigl(\Lambda_J(B)=0\bigr)\longrightarrow e^{-e^{-c}}.
\]
因此，内部坐标束屏幕的随机 exactification 具有一条 sharp 的 Poisson--Gumbel 临界窗，其极限分布完全由独立板层漏检的几何决定。
\end{theorem}
\begin{proof}
由定理 \ref{thm:conclusion-coordinate-bundle-random-audit-poisson-binomial}，\(\Lambda_J(B)\) 是各板层未命中事件之和。对均匀审计概率 \(p\)，每条板层的漏检概率都相同，且
\[
q_{m,s}(p)=\prod_{e\in E_{a_{J^c}}}(1-p)=(1-p)^{d_{m,s}}.
\]
板层总数为 \(L_J\)，故
\[
\Lambda_J(B)\sim \mathrm{Bin}\!\bigl(L_J,q_{m,s}(p)\bigr),
\]
并且
\[
\mathbb P\bigl(\Lambda_J(B)=0\bigr)=\bigl(1-q_{m,s}(p)\bigr)^{L_J}.
\]

若取
\[
p=p_{m,s}(c)=1-\exp\!\left(-\frac{\log L_J+c}{d_{m,s}}\right),
\]
则
\[
q_{m,s}(p)=\exp(-\log L_J-c)=\frac{e^{-c}}{L_J}.
\]
于是
\[
L_J q_{m,s}(p)\to e^{-c},
\qquad
q_{m,s}(p)\to 0.
\]
由经典二项分布 Poisson 极限定理，
\[
\Lambda_J(B)\Rightarrow \mathrm{Poisson}(e^{-c}).
\]
再对零点概率取极限便得到
\[
\mathbb P\bigl(\Lambda_J(B)=0\bigr)
\;=\;
\left(1-\frac{e^{-c}}{L_J}\right)^{L_J}
\longrightarrow
e^{-e^{-c}}.
\]
证毕。
\end{proof}

\begin{corollary}[临界审计密度的封闭式与 sharp 阈值]\label{cor:conclusion-coordinatebundle-critical-audit-density}
在定理 \ref{thm:conclusion-coordinatebundle-poisson-gumbel-critical-window} 的设定下，定义
\[
p_{m,s}^{\mathrm{crit}}:=p_{m,s}(0)=1-\exp\!\left(-\frac{\log L_J}{d_{m,s}}\right).
\]
则 exactification 的 sharp 阈值由单一无量纲参数
\[
\frac{\log L_J}{d_{m,s}}
=
\frac{m(n-s)\log 2}{2s\,2^{m(s-1)}}
\]
控制。更精确地：
\begin{enumerate}
  \item 若
  \[
  L_J(1-p)^{d_{m,s}}\to 0,
  \]
  则
  \[
  \mathbb P\bigl(\Lambda_J(B)=0\bigr)\to 1;
  \]
  \item 若
  \[
  L_J(1-p)^{d_{m,s}}\to\infty,
  \]
  则
  \[
  \mathbb P\bigl(\Lambda_J(B)=0\bigr)\to 0.
  \]
\end{enumerate}
此外，在
\[
\frac{\log L_J}{d_{m,s}}\to 0
\]
的 regime 中，成立 sharp 渐近
\[
p_{m,s}^{\mathrm{crit}}
\sim
\frac{\log L_J}{d_{m,s}}
=
\frac{m(n-s)\log 2}{2s\,2^{m(s-1)}}.
\]
特别地，对固定 \(s\ge 2\) 且 \(m\to\infty\)，临界审计密度按上式指数下沉。
\end{corollary}
\begin{proof}
第一条闭式来自
\[
L_J=2^{m(n-s)},
\qquad
d_{m,s}=2s\,2^{m(s-1)}.
\]
再由定理 \ref{thm:conclusion-coordinatebundle-poisson-gumbel-critical-window}，
\[
\mathbb P\bigl(\Lambda_J(B)=0\bigr)
=
\bigl(1-(1-p)^{d_{m,s}}\bigr)^{L_J}.
\]
若 \(L_J(1-p)^{d_{m,s}}\to 0\)，则右端趋于 \(1\)；若该量趋于 \(+\infty\)，则右端趋于 \(0\)。这就给出 sharp 阈值的充要尺度。

最后，当 \(\log L_J/d_{m,s}\to 0\) 时，
\[
1-e^{-x}\sim x
\qquad
(x\downarrow 0),
\]
故
\[
p_{m,s}^{\mathrm{crit}}
=
1-\exp\!\left(-\frac{\log L_J}{d_{m,s}}\right)
\sim
\frac{\log L_J}{d_{m,s}}.
\]
将 \(L_J\) 与 \(d_{m,s}\) 的闭式代入即可。证毕。
\end{proof}

\begin{theorem}[同层 Gibbs 占用、有效电阻与 Shapley 原子的严格同一]\label{thm:conclusion-coordinatebundle-resistance-shapley-identity}
\leanverified{paper\_conclusion\_coordinatebundle\_resistance\_shapley\_identity}
固定一条板层 \(E_{a_{J^c}}\)。在极小精确化 Gibbs 系综
\[
\mu_{J,\mathbf x}(B)\propto \prod_{e\in B}x_e
\]
下，对任意 \(e\in E_{a_{J^c}}\) 都有
\[
\mu_{J,\mathbf x}(e\in B)
=
\frac{x_e}{\sum_{f\in E_{a_{J^c}}}x_f}
=
x_e\,
R^{\mathrm{eff}}_{Q_{S_J^{\mathrm{int}}},\mathbf x}(\infty,a_{J^c}),
\]
其中
\[
R^{\mathrm{eff}}_{Q_{S_J^{\mathrm{int}}},\mathbf x}(\infty,a_{J^c})
=
\left(\sum_{f\in E_{a_{J^c}}}x_f\right)^{-1}.
\]

若进一步取无权情形 \(x_e\equiv 1\)，则对任意 \(e\in E_{a_{J^c}}\) 都有严格三重同一性
\[
\mu_{J,\mathbf 1}(e\in B)
=
R^{\mathrm{eff}}_{Q_{S_J^{\mathrm{int}}},\mathbf 1}(\infty,a_{J^c})
=
\operatorname{Sh}_e(g_J)
=
\frac{1}{|E_{a_{J^c}}|}.
\]
并且未归一化 Banzhaf 值满足
\[
\beta_e(g_J)=2^{1-|E_{a_{J^c}}|}.
\]
因此，同一板层上的物理边占用概率、电网络有效电阻与合作博弈中的 Shapley 原子在无权情形下并非仅仅同尺度，而是同一个数值对象。
\end{theorem}
\begin{proof}
由推论 \ref{cor:conclusion-coordinate-bundle-gibbs-resistance-closed-form}，
\[
\mu_{J,\mathbf x}(e\in B)=\frac{x_e}{\sum_{f\in E_{a_{J^c}}}x_f},
\qquad
R^{\mathrm{eff}}_{Q_{S_J^{\mathrm{int}}},\mathbf x}(\infty,a_{J^c})
=
\left(\sum_{f\in E_{a_{J^c}}}x_f\right)^{-1}.
\]
两式相乘即得加权恒等式。

在无权情形 \(x_e\equiv 1\) 下，
\[
\mu_{J,\mathbf 1}(e\in B)
=
R^{\mathrm{eff}}_{Q_{S_J^{\mathrm{int}}},\mathbf 1}(\infty,a_{J^c})
=
\frac{1}{|E_{a_{J^c}}|}.
\]
另一方面，推论 \ref{cor:conclusion-coordinate-bundle-face-shapley-banzhaf} 已给出
\[
\operatorname{Sh}_e(g_J)=\frac{1}{|E_{a_{J^c}}|},
\qquad
\beta_e(g_J)=2^{1-|E_{a_{J^c}}|}.
\]
故三重同一与 Banzhaf 公式同时成立。证毕。
\end{proof}

\begin{theorem}[严格边界价差下的纯星零温塌缩与熵密度]\label{thm:conclusion-coordinatebundle-pure-star-zero-temperature-entropy}
设 \(J\neq \varnothing\)，并在压缩图 \(\mathcal G_J=Q_{S_J^{\mathrm{int}}}\) 上赋边权。若满足严格价差
\[
\max\{w_e:\ e\ \text{incident to }\infty\}
<
\min\{w_e:\ e\ \text{not incident to }\infty\},
\]
则每一棵最小生成树都纯由外边界边组成；等价地，所有零温成本最优的极小精确化都属于纯外边界星树子族 \(\mathfrak E_J^{\min}\)。

若进一步取与上述 strict-gap 相容的对称特化，即所有外边界边同权且严格低于所有内部跨分量边，则零温最优族恰等于 \(\mathfrak E_J^{\min}\)，并且其精确基数为
\[
|\mathfrak E_J^{\min}|
=
\bigl(2s\,2^{m(s-1)}\bigr)^{L_J}
=
\bigl(2s\,2^{m(s-1)}\bigr)^{2^{m(n-s)}}.
\]
因此每个隐藏原子的零温熵密度恰为
\[
\frac{1}{L_J}\log |\mathfrak E_J^{\min}|
=
\log(2s)+m(s-1)\log 2.
\]
于是，一旦边界价差跨过 strict gap，内部坐标束屏幕的零温几何便塌缩为“每层各取一边”的纯星热力学，而其残余熵密度可由单层外边界的基数显式读出。
\end{theorem}
\begin{proof}
定理 \ref{thm:conclusion-coordinatebundle-strict-boundary-gap-forces-stars} 已证明：在 strict-gap 假设下，\(\mathcal G_J\) 的每一棵最小生成树都只能纯由 incident to \(\infty\) 的外边界边组成。由定理 \ref{thm:conclusion-coordinatebundle-all-minimal-exactifications-spanning-trees}，这等价于所有零温成本最优极小精确化都落在 \(\mathfrak E_J^{\min}\) 中。

在进一步的对称特化中，每层所有外边界边都同时达到同一最小代价，而所有内部跨分量边都被 strict gap 排除，故零温最优族恰好就是 \(\mathfrak E_J^{\min}\) 本身。再由推论 \ref{cor:conclusion-coordinatebundle-pure-boundary-star-family-count}，
\[
|\mathfrak E_J^{\min}|
=
\bigl(2s\,2^{m(s-1)}\bigr)^{L_J}.
\]
两边取对数并除以
\[
L_J=2^{m(n-s)}
\]
即得
\[
\frac{1}{L_J}\log |\mathfrak E_J^{\min}|
=
\log(2s)+m(s-1)\log 2.
\]
证毕。
\end{proof}

\begin{conjecture}[完全可积屏幕的坐标束刻画]\label{conj:conclusion-coordinatebundle-complete-integrability-characterization}
设 \(S\) 是某个固定 dyadic 复形上的部分屏幕，并同时满足：
\begin{enumerate}
  \item 对每个综合征 \(y\)，后验纤维 \(X\mid (Y_S=y)\) 都是仿射超立方体，且自由坐标相互独立；
  \item 在均匀外边界审计下，残余秩服从一条 sharp 的 Poisson--Gumbel 临界窗；
  \item 在无权零温极限中，同层 Gibbs 边占用概率、有效电阻与 Shapley 价值严格重合；
  \item 隐藏体积函数存在且仅存在一个能够消去全部高阶 Möbius 原子的标度。
\end{enumerate}
则在坐标重标与显然的可见等价之下，\(S\) 应当必为某个内部坐标束屏幕 \(S_J^{\mathrm{int}}\)。

这一猜想把猜想 \ref{conj:conclusion-coordinate-bundle-complete-factorization-rigidity} 的三条弱判据提升为概率、博弈、电网络与热力学的联合刚性：一旦后验超立方体、exactification 可乘性与 Möbius 单层支撑继续强化为 Poisson--Gumbel 窗口、对数唯一无交互标度与电阻--Shapley 同一，内部坐标束屏幕应当成为当前部分屏幕理论中唯一的完全可积分支。
\end{conjecture}

\noindent
因此，内部坐标束屏幕的隐藏几何可被进一步压成一条统一的可积链：原始尺度上的 Shannon 熵、exact audit 位数、线性补全维数与条件 Kolmogorov 容量都锁定在同一整数 \(L_J\) 上；而一旦切换到唯一无交互的对数标度，全部隐藏复杂度又严格分裂为 \(n-s\) 个 \(m\)-比特坐标原子。与此同时，均匀外边界审计的残余秩不只是一条 Poisson--binomial 和，还在临界窗内呈现出精确的 Poisson--Gumbel 截断；单层 Gibbs 边占用、电阻与 Shapley 值则在无权极限下三重合一，Banzhaf 指数只是同层基数的指数影子。最后，strict gap 把全部零温最优补全压缩为纯星子族，并留下每个隐藏原子都可显式读取的线性熵密度。由此，内部坐标束屏幕不再只是一般屏幕理论中的一个可解例子，而成为一个在后验几何、随机审计、合作博弈与电网络四个接口上同时闭合的完全可积对象。

\endinput
